\documentclass[aspectratio=169,11pt]{beamer}
\usepackage{paperdeck}
\title[GLU Variants]{GLU Variants Improve Transformer}
\subtitle{Noam Shazeer (2020), arXiv:2002.05202}
\author[Paper summary]{Paper summary for the CS-781 capstone}
\date{}
\begin{document}

\begin{frame}\titlepage\end{frame}

\paperglance
{N.\ Shazeer. \emph{GLU Variants Improve Transformer}. arXiv preprint, February 2020. Google.}
{Replace the Transformer's ReLU feed-forward layer with a \textbf{gated} version, such as SwiGLU or GEGLU, at equal parameters and compute. Perplexity and downstream scores improve.}
{Our A4 uses exactly the paper's parameter-matched SwiGLU: $d_{\text{model}}=768$, $d_{ff}=2048$, three matrices, no biases.}
{Main deck: ``A4: efficient decoder''. Glossary: ``Activation functions'' (SwiGLU as a valve).}

\begin{frame}{Roadmap}
  \begin{multicols}{2}\small\tableofcontents\end{multicols}
\end{frame}

% ------------------------------------------------------------------
\section{Background}

\begin{frame}{The feed-forward network (FFN) in a Transformer}
  \begin{columns}[T]
    \begin{column}{0.52\textwidth}
      \small
      Each Transformer block has attention (tokens exchange information) and an \textbf{FFN}, applied to each token on its own.
      \vspace{1mm}

      Original (Vaswani et al., 2017), and the bias-free T5 version used here:
      \[
        \text{FFN}_{\text{ReLU}}(x)=\max(xW_1,0)\,W_2
      \]
      Other activations had been proposed:
      \[
        \text{GELU}(x)=x\,\Phi(x),\qquad \text{Swish}_\beta(x)=x\,\sigma(\beta x)
      \]
      About two thirds of a Transformer block's parameters sit in the FFN, so its design matters.
    \end{column}
    \begin{column}{0.44\textwidth}
      \begin{exampleblock}{Intuition}
        \small
        Attention gathers facts from other tokens, e.g.\ ``he'' refers to ``respondent''.
        The FFN is where each token \emph{processes} what it gathered: a learned bank of feature detectors.
      \end{exampleblock}
      \begin{block}{Symbols}
        \small
        $x\in\mathbb{R}^{d}$: one token's vector. $W_1\in\mathbb{R}^{d\times d_{ff}}$, $W_2\in\mathbb{R}^{d_{ff}\times d}$. Typically $d_{ff}=4d$.
      \end{block}
    \end{column}
  \end{columns}
\end{frame}

\begin{frame}{Gated Linear Units (Dauphin et al., 2016)}
  \begin{columns}[T]
    \begin{column}{0.52\textwidth}
      \small
      A \textbf{GLU} multiplies two linear projections element-wise; one passes through a sigmoid:
      \[
        \text{GLU}(x)=\sigma(xW+b)\otimes(xV+c)
      \]
      \begin{itemize}
        \item $xV$: the \emph{content}.
        \item $\sigma(xW)$: a \emph{gate} between 0 and 1 deciding how much content passes.
      \end{itemize}
      Dropping the sigmoid gives the \textbf{bilinear} layer:
      \[
        \text{Bilinear}(x)=(xW+b)\otimes(xV+c)
      \]
    \end{column}
    \begin{column}{0.44\textwidth}
      \begin{exampleblock}{Analogy: a valve on a pipe}
        \small
        Water (content) flows through a pipe; a valve (gate) opens or closes.
        Both the water and the valve setting depend on the input.
        A plain ReLU is a one-way valve that can only be fully shut (negative) or open in proportion (positive).
      \end{exampleblock}
    \end{column}
  \end{columns}
\end{frame}

% ------------------------------------------------------------------
\section{The basic idea: GLU variants in the FFN}

\begin{frame}{The proposed variants}
  \begin{columns}[T]
    \begin{column}{0.54\textwidth}
      \small
      Replace the first linear layer and activation of the FFN with a GLU variant (no biases):
      \begin{align*}
        \text{FFN}_{\text{GLU}}(x)&=(\sigma(xW)\otimes xV)\,W_2\\
        \text{FFN}_{\text{Bilinear}}(x)&=(xW\otimes xV)\,W_2\\
        \text{FFN}_{\text{ReGLU}}(x)&=(\max(0,xW)\otimes xV)\,W_2\\
        \text{FFN}_{\text{GEGLU}}(x)&=(\text{GELU}(xW)\otimes xV)\,W_2\\
        \text{FFN}_{\text{SwiGLU}}(x)&=(\text{Swish}_1(xW)\otimes xV)\,W_2
      \end{align*}
      $\text{Swish}_1(z)=z\,\sigma(z)$, also called SiLU.
    \end{column}
    \begin{column}{0.42\textwidth}
      \begin{block}{Fair comparison: equal parameters}
        \small
        GLU variants have \textbf{three} matrices ($W$, $V$, $W_2$) instead of two.
        To keep parameters and compute equal, shrink $d_{ff}$ by $\tfrac23$:
        \[
          2\cdot d\cdot 3072 \;=\; 3\cdot d\cdot 2048
        \]
        For $d=768$: $4.72$M parameters per FFN in both cases.
      \end{block}
      \begin{exampleblock}{Our A4}
        \small
        \code{SwiGLU(d\_model=768, d\_ff=2048)}: identical to the paper's matched setting.
      \end{exampleblock}
    \end{column}
  \end{columns}
\end{frame}

\begin{frame}{Easy worked example (2 hidden units)}
  \begin{columns}[T]
    \begin{column}{0.5\textwidth}
      \small
      Suppose for one token: gate pre-activation $xW=[2,\,-1]$ and content $xV=[3,\,4]$.
      \vspace{1mm}

      \scriptsize
      \begin{tabular}{@{}ll@{}}
        \toprule
        Layer & Hidden vector (before $W_2$)\\
        \midrule
        ReLU (only $xW$) & $[2,\,0]$\\
        GLU & $\sigma(xW)\otimes xV=[2.64,\,1.08]$\\
        Bilinear & $xW\otimes xV=[6,\,-4]$\\
        ReGLU & $\max(0,xW)\otimes xV=[6,\,0]$\\
        SwiGLU & $\text{Swish}(xW)\otimes xV=[5.28,\,-1.08]$\\
        \bottomrule
      \end{tabular}
      \vspace{1mm}

      {\scriptsize $\sigma(2)=0.88$, $\sigma(-1)=0.27$; Swish$(2)=2\sigma(2)=1.76$, Swish$(-1)=-0.27$.}
    \end{column}
    \begin{column}{0.46\textwidth}
      \begin{exampleblock}{What to notice}
        \small
        \begin{itemize}
          \item In ReLU, unit 2 is simply off.
          \item In gated layers, the output is a \emph{product} of two input-dependent signals. The network can express ``this feature matters only when that one is present''.
          \item SwiGLU lets a small negative value through (smooth), unlike ReGLU's hard zero.
        \end{itemize}
      \end{exampleblock}
    \end{column}
  \end{columns}
\end{frame}

% ------------------------------------------------------------------
\section{Experiments}

\begin{frame}{Setup}
  \begin{columns}[T]
    \begin{column}{0.52\textwidth}
      \small
      \begin{itemize}
        \item Same code, model and task as the T5 base model: encoder--decoder, 12+12 layers, $d_{\text{model}}=768$, 12 heads, $d_k=d_v=64$.
        \item FFN: $d_{ff}=3072$ for ReLU, GELU and Swish; $d_{ff}=2048$ for the GLU variants.
        \item Pre-training: span-filling on C4, 524,288 steps, 128 examples $\times$ 512 input tokens per batch; Adafactor, inverse-square-root schedule, linear decay over the last 10\%.
        \item \textbf{No dropout during pre-training}, a departure from T5 that gave better results.
        \item Four short runs (65,536 steps) per variant measure run-to-run noise.
      \end{itemize}
    \end{column}
    \begin{column}{0.44\textwidth}
      \begin{block}{Measures}
        \small
        \begin{itemize}
          \item Held-out \textbf{log-perplexity} on the span-filling task (C4). Lower is better.
          \item Fine-tuning on a mixture of GLUE, SuperGLUE and SQuAD; development-set scores.
        \end{itemize}
      \end{block}
      \begin{exampleblock}{Cost}
        \small
        Each training step took about 0.15 s on a 32-core TPUv2 cluster. All variants cost the same by design.
      \end{exampleblock}
    \end{column}
  \end{columns}
\end{frame}

\begin{frame}{Pre-training perplexity}
  \begin{columns}[T]
    \begin{column}{0.52\textwidth}
      \footnotesize
      Held-out log-perplexity (lower is better); standard deviation over 4 runs in brackets.
      \vspace{1mm}

      \begin{tabular}{@{}lcc@{}}
        \toprule
        FFN variant & 65,536 steps & 524,288 steps\\
        \midrule
        ReLU (baseline) & 1.997 (0.005) & 1.677\\
        GELU & 1.983 (0.005) & 1.679\\
        Swish & 1.994 (0.003) & 1.683\\
        GLU & 1.982 (0.006) & 1.663\\
        Bilinear & 1.960 (0.005) & 1.648\\
        \textbf{GEGLU} & \textbf{1.942} (0.004) & \textbf{1.633}\\
        \textbf{SwiGLU} & 1.944 (0.010) & 1.636\\
        ReGLU & 1.953 (0.003) & 1.645\\
        \bottomrule
      \end{tabular}
    \end{column}
    \begin{column}{0.44\textwidth}
      \begin{exampleblock}{Reading it}
        \small
        \begin{itemize}
          \item Changing only the activation (GELU, Swish) barely helps.
          \item \textbf{Every gated variant beats every ungated one.}
          \item GEGLU and SwiGLU are best: about 0.04 lower than ReLU after full training, far above the run-to-run noise of about 0.005.
          \item Even the bilinear layer (no non-linearity in the gate) beats ReLU.
        \end{itemize}
      \end{exampleblock}
    \end{column}
  \end{columns}
\end{frame}

\begin{frame}{Downstream fine-tuning (development sets)}
  \begin{columns}[T]
    \begin{column}{0.52\textwidth}
      \scriptsize
      \begin{tabular}{@{}lccc@{}}
        \toprule
        FFN variant & GLUE & SGLUE & SQuAD F1\\
        \midrule
        ReLU & 83.80 & 72.76 & 90.87\\
        GELU & 83.86 & 72.98 & 90.79\\
        Swish & 83.60 & 72.40 & 90.76\\
        GLU & 84.20 & 73.95 & 90.69\\
        GEGLU & 84.12 & 73.96 & 91.12\\
        Bilinear & 83.79 & 73.81 & 91.06\\
        SwiGLU & 84.36 & \textbf{74.56} & 91.03\\
        ReGLU & \textbf{84.67} & 73.66 & \textbf{91.18}\\
        \midrule
        T5 paper baseline & 83.28 & 71.36 & 88.81\\
        \bottomrule
      \end{tabular}
      \vspace{1mm}

      {\scriptsize T5's reported standard deviations: GLUE 0.235, SuperGLUE 0.416.}
    \end{column}
    \begin{column}{0.44\textwidth}
      \begin{exampleblock}{Reading it}
        \small
        \begin{itemize}
          \item GLU variants are best on most tasks, e.g.\ SuperGLUE 74.56 (SwiGLU) vs 72.76 (ReLU).
          \item The results are noisy: each model was fine-tuned only once.
          \item The paper's ReLU model beats the original T5 numbers (83.80 vs 83.28). The author attributes this to removing dropout during pre-training.
        \end{itemize}
      \end{exampleblock}
    \end{column}
  \end{columns}
\end{frame}

% ------------------------------------------------------------------
\section{Assessment and outlook}

\begin{frame}{Relevance to our capstone}
  \begin{columns}[T]
    \begin{column}{0.5\textwidth}
      \small
      \begin{itemize}
        \item A4 uses SwiGLU; A1 and A3 use GELU; A2 uses ReLU. This makes the FFN one of the \emph{differences} between our architectures.
        \item The paper predicts a small but real perplexity gain for A4 from the FFN alone.
        \item Our default \code{dropout: 0.0} matches the paper's finding that no dropout helps pre-training when data is plentiful.
      \end{itemize}
    \end{column}
    \begin{column}{0.46\textwidth}
      \begin{block}{A controlled ablation you could run}
        \small
        Train A1 twice, identical except the FFN:
        \begin{itemize}
          \item GELU, $d_{ff}=3072$
          \item SwiGLU, $d_{ff}=2048$
        \end{itemize}
        Same parameters, same data. Compare validation loss and dev BERTScore.
        This separates the FFN effect from A4's other changes (GQA, RoPE, sliding window).
      \end{block}
    \end{column}
  \end{columns}
\end{frame}

\begin{frame}{Strengths and limitations}
  \begin{columns}[T]
    \begin{column}{0.48\textwidth}
      \begin{exampleblock}{Strengths}
        \small
        \begin{itemize}
          \item A tiny code change with consistent gains.
          \item A fair comparison: parameters and compute matched.
          \item Measures run-to-run noise for perplexity.
          \item No extra cost at training or inference.
          \item Adopted widely in later large models.
        \end{itemize}
      \end{exampleblock}
    \end{column}
    \begin{column}{0.48\textwidth}
      \begin{alertblock}{Limitations}
        \small
        \begin{itemize}
          \item \textbf{No explanation} of why gating works. The author jokingly credits ``divine benevolence''.
          \item One model size and one architecture (T5 base).
          \item Fine-tuning results come from one run each, and are noisy.
          \item Pre-training changes (no dropout) are mixed into the comparison with T5's numbers.
          \item A very short paper, with no ablation of $d_{ff}$ or scaling behaviour.
        \end{itemize}
      \end{alertblock}
    \end{column}
  \end{columns}
\end{frame}

\begin{frame}{Future research directions}
  \begin{columns}[T]
    \begin{column}{0.5\textwidth}
      \small
      \textbf{What followed}
      \begin{itemize}
        \item SwiGLU became standard in large language models, e.g.\ PaLM (Chowdhery et al., 2022) and LLaMA (Touvron et al., 2023).
        \item GEGLU is used in the T5 v1.1 checkpoints and in the Gemma models.
        \item Sparse FFNs: mixture-of-experts layers (e.g.\ Switch Transformer, Fedus et al., 2022) replace one FFN with many expert FFNs, and only a few run per token.
      \end{itemize}
    \end{column}
    \begin{column}{0.46\textwidth}
      \begin{block}{Open questions}
        \small
        \begin{itemize}
          \item \textbf{Why} does multiplicative gating help: better optimisation, or more expressive functions?
          \item Does the gain grow, shrink or vanish at larger scale?
          \item What is the best $d_{ff}$ ratio for gated FFNs?
          \item Do gated FFNs store factual knowledge differently?
          \item For legal NLP: do gated FFNs help with rare legal terms, or with long formulaic sentences?
        \end{itemize}
      \end{block}
    \end{column}
  \end{columns}
\end{frame}

\begin{frame}{Summary}
  \begin{itemize}
    \item The Transformer FFN is $\max(xW_1,0)W_2$. Changing only the activation (GELU, Swish) gives little.
    \item \textbf{Gated} FFNs multiply a content projection by an activated gate projection: GLU, Bilinear, ReGLU, GEGLU, SwiGLU.
    \item At equal parameters ($d_{ff}$ reduced by $\tfrac23$), all gated variants lower pre-training perplexity; GEGLU and SwiGLU are best (1.633 and 1.636 vs 1.677).
    \item Downstream scores improve on most tasks, e.g.\ SuperGLUE 74.56 vs 72.76.
    \item For us: the reason A4 uses SwiGLU with $d_{ff}=2048$.
  \end{itemize}
\end{frame}

\begin{frame}{Reference}
  \small
  Shazeer, N.\ (2020). GLU Variants Improve Transformer. arXiv:2002.05202.
  \vspace{3mm}

  {\footnotesize Other work cited: Dauphin et al.\ (2016) Language Modeling with Gated Convolutional Networks; Hendrycks and Gimpel (2016) GELU; Ramachandran et al.\ (2017) Swish; Raffel et al.\ (2020) T5; Chowdhery et al.\ (2022) PaLM; Fedus et al.\ (2022) Switch Transformers; Touvron et al.\ (2023) LLaMA.}
\end{frame}

\end{document}
